%=======================================================================
%  Queries, suspected misprints and points of clarification concerning
%  D. J. Steigmann and M. Shirani,
%  "Principles of Continuum Mechanics", World Scientific.
%
%  Compile:  latexmk -pdf book_questions.tex
%=======================================================================
\documentclass[12pt,a4paper]{article}

\usepackage[margin=2.2cm]{geometry}
\usepackage{amsmath,amssymb,amsthm}
\usepackage{accents}
\DeclareRobustCommand{\vek}[1]{\underaccent{\tilde}{#1}}
\usepackage{enumitem}
\usepackage[colorlinks=true,linkcolor=blue!55!black,hypertexnames=false]{hyperref}

\newcommand{\bv}{\vek{v}}
\newcommand{\bu}{\vek{u}}
\newcommand{\bx}{\vek{x}}
\newcommand{\bX}{\vek{X}}
\newcommand{\bn}{\vek{n}}
\newcommand{\bF}{\vek{F}}
\newcommand{\bC}{\vek{C}}
\newcommand{\bE}{\vek{E}}
\newcommand{\hbE}{\hat{\vek{E}}}
\newcommand{\be}{\vek{e}}
\newcommand{\bA}{\vek{A}}
\newcommand{\ba}{\vek{a}}
\newcommand{\bb}{\vek{b}}
\newcommand{\bze}{\vek{0}}
\newcommand{\bU}{\vek{U}}
\newcommand{\bM}{\vek{M}}
\newcommand{\bY}{\vek{Y}}
\newcommand{\bZ}{\vek{Z}}
\newcommand{\hI}{\hat{\vek{I}}}
\newcommand{\bB}{\vek{B}}
\newcommand{\bV}{\vek{V}}
\newcommand{\by}{\vek{y}}
\newcommand{\I}{\vek{I}}
\newcommand{\bchi}{\vek{\chi}}
\newcommand{\bkappa}{\kappa}
\newcommand{\dd}{\delta}
\newcommand{\bL}{\vek{L}}
\newcommand{\bw}{\vek{w}}
\newcommand{\bQ}{\vek{Q}}
\newcommand{\bc}{\vek{c}}
\newcommand{\bd}{\vek{d}}
\newcommand{\bT}{\vek{T}}
\newcommand{\bS}{\vek{S}}
\newcommand{\bP}{\vek{P}}
\newcommand{\bG}{\vek{G}}
\newcommand{\Et}{\mathbb{E}^{3}}
\newcommand{\Vt}{\mathbb{E}^{3}}
\newcommand{\grad}{\operatorname{grad}}
\newcommand{\tr}{\operatorname{tr}}

\newenvironment{query}[1]%
  {\par\medskip\noindent\textbf{#1}\quad}{\par\medskip}

\title{Queries and suspected misprints\\
\large\textsc{Principles of Continuum Mechanics}\\
\large D. J. Steigmann and M. Shirani}
\author{Reader's notes}
\date{\today}

\begin{document}
\maketitle

\noindent
The following notes were accumulated during a line-by-line reading of
Chapters 1--12. They are grouped by confidence: Part I contains items I
believe to be genuine misprints, verified by independent computation;
Part II contains statements that appear to be stronger or looser than
intended; Part III contains notational and pedagogical points where a
sentence of clarification would help the reader; Part IV concerns the
figures. Equation numbers and page numbers refer to the printed text.

%=======================================================================
\section*{Part I. Verified misprints}
%=======================================================================

\begin{query}{1. Equation (2.32), p.~45.}
The printed spatial time derivatives are
\[
   \tilde v_2'=\frac{1-2t^{2}}{1+t^{2}}\,x_2-\frac{2t}{1+t^{2}}\,x_3,
   \qquad
   \tilde v_3'=\frac{1-2t^{2}}{1+t^{2}}\,x_3-\frac{2t}{1+t^{2}}\,x_2 .
\]
From the inversion (2.31) the spatial velocity components are
$\tilde v_2=(1+t^{2})^{-1}(tx_2+x_3)$ and
$\tilde v_3=(1+t^{2})^{-1}(tx_3-x_2)$, whence, by the quotient rule at
fixed $\bx$,
\[
   \tilde v_2'=\frac{(1+t^{2})-t(2t)}{(1+t^{2})^{2}}\,x_2
              -\frac{2t}{(1+t^{2})^{2}}\,x_3
             =\frac{1-t^{2}}{(1+t^{2})^{2}}\,x_2-\frac{2t}{(1+t^{2})^{2}}\,x_3,
\]
and similarly for $\tilde v_3'$. Both the numerator ($1-t^{2}$, not
$1-2t^{2}$) and the exponent of the denominator differ from the printed
values.

That the corrected expressions are the right ones is confirmed by the
consistency check the text itself proposes. Using (2.29) with the printed
matrix (2.33), which is correct,
\[
\begin{aligned}
   \tilde a_2&=\tilde v_2'+\tilde v_{2,2}\tilde v_2+\tilde v_{2,3}\tilde v_3\\
   &=\frac{(1-t^{2})x_2-2tx_3}{(1+t^{2})^{2}}
     +\frac{t}{1+t^{2}}\cdot\frac{tx_2+x_3}{1+t^{2}}
     +\frac{1}{1+t^{2}}\cdot\frac{tx_3-x_2}{1+t^{2}}\\
   &=\frac{(1-t^{2}+t^{2}-1)x_2+(-2t+t+t)x_3}{(1+t^{2})^{2}}=0,
\end{aligned}
\]
in agreement with the stated result $\tilde a_2=0$. With the printed form
of (2.32) the same computation does not vanish, so the misprint is
detectable from within the example.
\end{query}

\begin{query}{2. Equation (3.6), p.~64.}
The conjugate components are printed as
\[
   \bar{\bu}=\bar u_j\be_j \quad\text{with}\quad u_j=a_j-\mathrm{i}b_j .
\]
Since (3.4) already sets $u_j=a_j+\mathrm{i}b_j$, the left-hand side of the
second relation should read $\bar u_j$. As printed, (3.4) and (3.6)
together assert $a_j+\mathrm{i}b_j=a_j-\mathrm{i}b_j$.
\end{query}

\begin{query}{3. Problem 9 of Chapter 3, p.~89.}
The problem asks the reader to show that, for the two-dimensional simple
shear $\bF=\vek{1}+\gamma\,\be_{1}\otimes\hbE_{2}$,
\[
   \bU=\frac{1}{\sqrt{4+\gamma^{2}}}
   \Big[2\hI+\gamma(\hbE_{1}\otimes\hbE_{2}+\hbE_{2}\otimes\hbE_{1})
   +\gamma^{2}\,\hbE_{1}\otimes\hbE_{1}\Big].
\]
The last dyad should be $\hbE_{2}\otimes\hbE_{2}$. With the printed
version, in matrix form
$\bU=\tfrac1I\left(\begin{smallmatrix}2+\gamma^{2}&\gamma\\
\gamma&2\end{smallmatrix}\right)$ with $I=\sqrt{4+\gamma^{2}}$, one gets
\[
   \bU^{2}=\begin{pmatrix}1+\gamma^{2}&\gamma\\ \gamma&1\end{pmatrix},
\]
which is the matrix of $\bB$, not of $\bC$; the required
$\bU^{2}=\bC=\left(\begin{smallmatrix}1&\gamma\\ \gamma&1+\gamma^{2}
\end{smallmatrix}\right)$ of $(3.55)$ is obtained only with
$\gamma^{2}\hbE_{2}\otimes\hbE_{2}$:
\[
   \bU^{2}=\frac{1}{4+\gamma^{2}}
   \begin{pmatrix}4+\gamma^{2}&\gamma(4+\gamma^{2})\\
   \gamma(4+\gamma^{2})&(4+\gamma^{2})(1+\gamma^{2})\end{pmatrix}
   =\begin{pmatrix}1&\gamma\\ \gamma&1+\gamma^{2}\end{pmatrix},
\]
the last entry using
$\gamma^{2}+(2+\gamma^{2})^{2}=4+5\gamma^{2}+\gamma^{4}
=(4+\gamma^{2})(1+\gamma^{2})$.

The printed expression is in fact the correct formula for the \emph{left}
stretch $\bV$, with hats attached: applying the same identity
$I\bV=\bB+J\I$ to $(3.56)$ gives
\[
   \bV=\frac{1}{\sqrt{4+\gamma^{2}}}
   \Big[2\I+\gamma(\be_{1}\otimes\be_{2}+\be_{2}\otimes\be_{1})
   +\gamma^{2}\,\be_{1}\otimes\be_{1}\Big],
\]
so the two seem to have been interchanged. The distinction is exactly the
one the text itself stresses between $(3.53)$ and $(3.54)$, where the
coefficient $1+\gamma^{2}$ sits on $\hbE_{2}\otimes\hbE_{2}$ in $\bC$ and
on $\be_{1}\otimes\be_{1}$ in $\bB$.
\end{query}

\begin{query}{4. The component form of $\bL$ in the spinning disc, p.~96.}
Immediately below (4.27) the text writes ``Thus,
$L_{ij}=\omega e_{ijk}n_{j}$'' and, two lines further on,
``$W_{ij}=\omega e_{ijk}n_{j}$''. Neither can stand as printed: the index
$j$ appears three times, once free on the left and twice on the right, so
the equations are not index balanced, and $k$ is left free on the right
with no counterpart on the left.

Equation (4.27) itself,
\[
   v_{i}(\by,t)-v_{i}(\bx,t)=\omega e_{ijk}n_{j}(y_{k}-x_{k}),
\]
is correct. Comparing it with $v_{i}(\by,t)-v_{i}(\bx,t)
=L_{ik}(y_{k}-x_{k})$ identifies $i$ and $k$ as the free indices and $j$
as the summed one, so the two lines should read
\[
   L_{ik}=\omega e_{ijk}n_{j},
   \qquad
   W_{ik}=\omega e_{ijk}n_{j}.
\]
The intervening conclusion is then immediate,
$L_{ki}=\omega e_{kji}n_{j}=-\omega e_{ijk}n_{j}=-L_{ik}$, so $D_{ik}=0$
and the motion is rigid, as the text says. Nothing downstream is affected.

A second, smaller slip occurs in the chain (4.28). Its second line is
printed as $\tfrac12\omega e_{jki}L_{ik}\be_{j}$ and its third as
$\tfrac12\omega e_{jki}e_{lki}n_{l}\be_{j}$, so the factor $\omega$
appears in both. Since $L_{ik}$ already carries $\omega$, the factor
should enter exactly once, at the step where $L$ is substituted:
\[
   \bw=\tfrac12 e_{jki}v_{i,k}\be_{j}=\tfrac12 e_{jki}L_{ik}\be_{j}
   =\tfrac12\,\omega\,e_{jki}e_{ilk}n_{l}\be_{j}
   =\tfrac12\,\omega\,(2\delta_{jl})n_{l}\be_{j}=\omega\bn .
\]
The stated result $\bw=\omega\bn$ is of course correct.
\end{query}

\begin{query}{5. A sign in equation (6.116), p.~134.}
The printed equation reads
\[
   e_{ijk}x_{j}\big\{\rho(\dot v_{k}-b_{k})+T_{kl,l}\big\}=e_{ijk}T_{kj}.
\]
The sign inside the braces should be negative. Expanding
$\rho\dot\psi_{i}=\rho r_{i}+A_{il,l}$ with the identifications (6.114),
\[
   \rho e_{ijk}x_{j}\dot v_{k}
   =\rho e_{ijk}x_{j}b_{k}+e_{ijk}T_{kj}+e_{ijk}x_{j}T_{kl,l},
\]
and collecting gives $e_{ijk}x_{j}\{\rho(\dot v_{k}-b_{k})-T_{kl,l}\}
=e_{ijk}T_{kj}$.

The sentence immediately after the equation requires the minus sign as
well: it says the braces vanish by (6.110). Since (6.110) reads
$\rho\dot v_{k}=\rho b_{k}+T_{kl,l}$, the combination that vanishes is
$\rho(\dot v_{k}-b_{k})-T_{kl,l}$. With the printed plus sign the braces
would equal $2T_{kl,l}$ on solutions of (6.110), and the symmetry of the
Cauchy stress would not follow.
\end{query}

\begin{query}{6. A missing transpose in equation (9.14), p.~183.}
The second term of the printed (9.14) is
$\dot\bQ_{1}\bQ^{t}_{1}(\bc_{1}-\bc_{2})$. It should be
$\dot\bQ^{t}_{1}(\bc_{1}-\bc_{2})$.

The error is visible without computing anything, from the spaces the
factors live in. The vectors $\bc_{k}$ lie in $\mathbb{E}^{3+}$, since
(9.6) adds them to $\bQ_{k}\bx_{k}$ there. Now $\bQ^{t}_{1}$ carries
$\mathbb{E}^{3+}$ to $\Et$ and $\dot\bQ_{1}$ carries $\Et$ back to
$\mathbb{E}^{3+}$, so the printed expression lands in
$\mathbb{E}^{3+}$. But the vector it is added to must lie in $\Et$,
because (9.13) adds it to $\bQ\bv_{1}$ there. Exactly one transpose is
missing.

A check of the corrected form: with $\bQ_{2}=\bQ_{1}$ one has
$\bQ=\I$ and (9.12) reduces to a pure translation, for which the
corrected term reproduces the relative velocity correctly while the
printed one does not.
\end{query}

\begin{query}{7. A repeated index in equation (9.35), p.~191.}
The determinant is printed as
\[
   (\eta_{1}-\eta_{1})(\eta_{2}-\eta_{3})(\eta_{3}-\eta_{1}),
\]
whose first factor is identically zero. The intended factor is
$\eta_{1}-\eta_{2}$, giving the Vandermonde form
$(\eta_{1}-\eta_{2})(\eta_{2}-\eta_{3})(\eta_{3}-\eta_{1})$. The
argument at that point requires the determinant to be non-zero when the
three principal values are distinct, and with the printed factor it
vanishes always, which would destroy the conclusion.
\end{query}

\begin{query}{8. The $11$ component of $\bB^{2}$ in (11.37) and (11.39),
p.~213--214.}
Both displays print the $11$ component of $\bB^{2}$ in simple shear as
$(1+\gamma^{2})+\gamma^{2}$. It should be
$(1+\gamma^{2})^{2}+\gamma^{2}$. With
\[
   (B_{ij})=\begin{pmatrix}1+\gamma^{2}&\gamma&0\\
   \gamma&1&0\\ 0&0&1\end{pmatrix},
\]
squaring gives
$B^{2}_{11}=B_{11}B_{11}+B_{12}B_{21}=(1+\gamma^{2})^{2}+\gamma^{2}$. At
$\gamma=1$ the correct value is $5$ and the printed expression gives
$3$; direct squaring of
$\left(\begin{smallmatrix}2&1&0\\1&1&0\\0&0&1\end{smallmatrix}\right)$
confirms $5$. The other two entries of the printed $\bB^{2}$, namely
$1+\gamma^{2}$ and $\gamma(2+\gamma^{2})$, are correct.

The consequence is not confined to that entry. Equation (11.42),
$T_{11}-T_{22}=\gamma T_{12}$, is derived from these components and is
described in the text as a universal relation furnishing a critical test
of the theory. With the printed value it comes out as
$T_{11}-T_{22}=\gamma^{2}(f_{1}+f_{2})$, which is not $\gamma T_{12}$
unless $f_{2}=0$. With the corrected value it comes out exactly.
\end{query}

\begin{query}{9. A spurious factor $\gamma$ in equation (11.40),
p.~214.}
The printed shear response is
\[
   T_{12}=\gamma\mu(\gamma^{2}),\qquad
   \mu(\gamma^{2})=f_{1}(\gamma^{2})
   +\gamma\big(2+\gamma^{2}\big)f_{2}(\gamma^{2}).
\]
The factor $\gamma$ in the second term does not belong; the correct
expression is $\mu(\gamma^{2})=f_{1}(\gamma^{2})
+(2+\gamma^{2})f_{2}(\gamma^{2})$.

Two independent checks. First, the $12$ entry of (11.39) is
$f_{1}\gamma+f_{2}\gamma(2+\gamma^{2})$, and factoring out the $\gamma$
that (11.40) displays leaves $f_{1}+(2+\gamma^{2})f_{2}$. Second, the
notation $\mu(\gamma^{2})$ asserts that $\mu$ depends on $\gamma$ only
through $\gamma^{2}$, whereas the printed expression contains a term odd
in $\gamma$ and so contradicts its own argument list. The same spurious
factor also breaks (11.42), which is derived from $\mu$ two lines later.
\end{query}

%=======================================================================
\section*{Part II. Statements that appear stronger or looser than intended}
%=======================================================================

\begin{query}{10. The sign of $\det\bF$ and occupiability, p.~50.}
The text states: ``If $\kappa$ is not occupiable, then $\det\bF<0$ and
(2.50)$_2$ is replaced by $J=|\det\bF|$.''

A non-occupiable reference configuration does not by itself force
$\det\bF<0$. What the preceding argument establishes is the converse
direction: if $\kappa$ is occupiable then $\det\bF>0$. A reference
configuration may fail to be occupiable while still being reachable, at
each material point, by an orientation-preserving map; the obstruction to
occupiability is global, whereas the sign of $\det\bF$ is a pointwise
matter. Would ``$\det\bF$ need not be positive, and (2.50)$_2$ is then
replaced by $J=|\det\bF|$'' be the intended reading? If the stronger
statement is intended, a word on why non-occupiability forces the
orientation reversal would be very welcome, since this is the first place
in the book where the reader is asked to accept a sign convention on
faith.
\end{query}

\begin{query}{11. Equation (2.9) and the term ``material velocity'', p.~40.}
Equation (2.9) introduces
$\tilde\bv(\bx,t)=\bv(\bchi^{-1}(\bx,t),t)$ with the words
``Alternatively, the material velocity is given by the function
$\tilde\bv$''. In \S2.2, however, the tilde is reserved for the
\emph{spatial} description, and $\tilde\bv$ is precisely the spatial
description of the velocity field. Since \S2.2 goes to some length to
insist that $\bar f$, $\tilde f$, $\hat f$ are three different functions
with three different domains, calling $\tilde\bv$ ``the material
velocity'' one page earlier is likely to be read as a contradiction. Is
``the velocity in the spatial description'' intended? The same phrase
recurs after (2.17).
\end{query}

\begin{query}{12. Equation (3.19) and the labels $\lambda_1,\lambda_2$, p.~66.}
In \S3.2, $\lambda_1$ and $\lambda_2$ are defined as the minimum and
maximum of $f(\bu)=\bu\cdot\bA\bu$ on the unit sphere, and $\bu_1,\bu_2$
as the corresponding minimizers and maximizers. Everywhere else in the
book, including (3.15) and (3.38), the subscript on $\lambda_i$ is a
plain enumeration index over the three principal values. In a
three-dimensional problem the extremal pair is the smallest and the
largest, so the local convention of \S3.2 makes $\lambda_2$ denote the
\emph{largest} principal value while the global convention leaves it as
the middle one. Would $\lambda_{\min},\lambda_{\max}$ (or
$\lambda_-,\lambda_+$) in \S3.2 avoid the clash? Equation (3.26) is where
the two conventions are hardest to hold apart.
\end{query}

\begin{query}{13. Problem 21(g), Chapter 1, p.~35.}
The identity $\det(\bA+\ba\otimes\bb)=(\det\bA)(1+\ba\cdot\bA^{-t}\bb)$
presupposes that $\bA$ is invertible, but the hypothesis is not stated,
and parts (a)--(f) of the same problem carry no such restriction. Part (d)
of the problem does state its hypothesis implicitly by writing
$\bA^{-1}$; part (g) does so only through the $\bA^{-t}$ on the
right-hand side, which a student is likely to reach only after having
attempted the general case.
\end{query}

\begin{query}{14. Equation (3.53), third line, p.~73.}
The cross terms are printed as
$\gamma(\hbE_2\otimes\vek{1}^{t}\bE_1+\vek{1}^{t}\bE_1\otimes\hbE_2)$.
Expanding $(\vek{1}+\gamma\be_1\otimes\hbE_2)^{t}
(\vek{1}+\gamma\be_1\otimes\hbE_2)$ by the rules of Problem 17 produces
$\vek{1}^{t}\be_1$, not $\vek{1}^{t}\bE_1$: the vector being shifted back
is the \emph{spatial} basis vector $\be_1$ appearing in (3.51). The two
agree only after the identification $\{\bE_A\}=\{\be_i\}$ made at (3.49),
and the passage to the fourth line, $\vek{1}^{t}\bE_1=\hbE_1$, uses
(2.68) directly and therefore really does want $\bE_1$. Since \S2.6 is at
pains to explain that $\be_i$ and $\bE_A$ live in the same space but are
not the same vectors in general, could the identification be recalled
explicitly at this step? As it stands, the example silently uses a
convention adopted four equations earlier for an unrelated purpose.
\end{query}

\begin{query}{15. Equation (3.57), p.~74.}
The chain
$\det\bU=[\bU\bu_1,\bU\bu_2,\bU\bu_3]/[\bu_1,\bu_2,\bu_3]
=[\lambda_1\bu_1,\lambda_2\bu_2,\lambda_3\bu_3]=\lambda_1\lambda_2\lambda_3$
drops the denominator between the second and third members. The step is
of course legitimate, since $[\bu_1,\bu_2,\bu_3]=\pm1$ for an orthonormal
triad and the sign cancels against the numerator, but as printed the
equality signs do not hold term by term.
\end{query}

%=======================================================================
\section*{Part III. Notation and points where a clarifying sentence would help}
%=======================================================================

\begin{query}{16. The symbol $e$ for both basis vectors and the permutation
symbol, Chapter 1.}
In (1.164) the single line
\[
   \be_j\cdot\operatorname{curl}\bv=\operatorname{div}(v_i\be_i\times\be_j)
   =\operatorname{div}(v_i e_{ijk}\be_k)=(v_ie_{ijk})_{,k}=e_{kij}v_{i,k}
\]
uses $e$ as a basis vector and as the alternating symbol in the same
product. The distinction is carried entirely by the boldface, which at
the size used for subscripted symbols is easy to lose. Many texts use
$\varepsilon_{ijk}$ for the permutation symbol for exactly this reason.
Is there a reason for the present choice that I have missed? The same
collision occurs in Problem 11(a) of Chapter 1 and in Problem 1 of
Chapter 2.
\end{query}

\begin{query}{17. The symbol $\dd_{iA}$ in (2.70), p.~54.}
The shifter is decomposed as $\vek{1}=\dd_{iA}\be_i\otimes\hbE_A$, and the
text then observes that ``$\dd_{iA}$ are the components of the usual
Kronecker delta if, and only if, $\{\bE_A\}=\{\be_i\}$. Otherwise, they
are the cosines of the fixed angles between the elements of the two
sets.''

This is clear once read, but the notation works against the sentence: the
symbol $\dd$ is used throughout the book with the meaning
$\dd_{ij}=\be_i\cdot\be_j$, and (2.74) then reads
$\dd_{iA}\dd_{iB}=\dd_{AB}$, in which the two $\dd$'s on the left have the
general meaning and the one on the right has the Kronecker meaning. Would
a distinct symbol for the shifter components, $s_{iA}$ say, be worth the
extra letter? This is the one place in Chapter 2 where I had to reread
several times to be sure which object was which.
\end{query}

\begin{query}{18. Equation (2.12), p.~40: the absence of the subscript
$\kappa$.}
I record this as a query answered rather than as a correction, since it is
the first thing a reader checks. Equation (2.12) reads
$\tilde f(\bx,t)=\bar f(\bchi^{-1}(\bx,t),t)$ with no subscript on
$\bchi$, whereas the very next equation, (2.13), carries $\bchi_\kappa$
throughout, and the reader naturally suspects an omission. It is not one:
$\bchi^{-1}$ in (2.12) is the inverse of the motion (2.1), which returns
the \emph{material point} $p$, and $\bar f$ is precisely the function that
eats $p$, so the unsubscripted map is the correct one; $\bchi_\kappa^{-1}$
would return $\bX$ and would have to be composed with $\hat f$ instead.
The same applies to (2.9).

The point is worth a half-sentence in the text, because $\bchi$ and
$\bchi_\kappa$ differ by the fixed bijection $\kappa$ alone and the eye
does not distinguish them; the reader who assumes a misprint here will
misread the entire section. A parenthetical ``(note that (2.1) rather than
(2.6) is inverted here, since $\bar f$ takes the material point as its
argument)'' would settle it.
\end{query}

\begin{query}{19. $\bx'=\bze$ in the sanity check preceding \S2.4, p.~44.}
The line ``application of (2.20) to the position vector $\bx$, combined
with $\bx'=\bze$, yields
$\dot x_i=x_{i,j}\tilde v_j=\dd_{ij}\tilde v_j=\tilde v_i$, as expected''
is correct but is delivered without comment, and $\bx'=\bze$ reads as a
flat contradiction of the fact that material points move. The content is
that the field in question, taken in the spatial description, is
$\tilde\bu(\bx,t)=\bx$, the identity map of $\mathcal{E}^{3}$, which
carries no $t$ at all; the motion resides entirely in $\hat\bu=\bchi_\kappa$.
Given how much care \S2.2 devotes to the distinction between $'$ and
$\dot{\phantom{f}}$, and given that this is the extreme illustration of it
($\bze$ against $\bv$ for the same field), one line of explanation here
would repay itself. As printed, ``as expected'' invites the reader to
believe the step is routine when in fact it is the sharpest test of
whether \S2.2 has been understood.
\end{query}

\begin{query}{20. Section title of \S2.2 versus the order in (2.11).}
The section is titled ``Material, referential and spatial descriptions'',
whereas (2.11) lists $\bar f$, $\tilde f$, $\hat f$ and the sentence after
(2.13) reads ``the material, spatial and referential descriptions,
respectively''. The orders differ. A trivial matter, but the reader who
uses the section title as a mnemonic for which decoration is which will be
misled.
\end{query}

\begin{query}{21. Equation (2.96), p.~58.}
The step $\tfrac12(\bM\cdot\bC\bM-1)=\bM\cdot\bE\bM$ uses
$1=\bM\cdot\hat{\I}\bM$, which holds because $|\bM|=1$. The unit-length
hypothesis is stated at (2.66) but is not recalled here, and the reader
tracking the algebra will look for a missing factor. Similarly at (2.98).
\end{query}

%=======================================================================
\section*{Part IV. Figures}
%=======================================================================

\begin{query}{22. Figure 2.5, p.~41.}
The caption reads ``Different material points occupy position $\bx$ at
different times'', but the figure shows two disjoint regions. Since the
whole point of the spatial description is that \emph{one} body sweeps past
\emph{one} fixed place, a single body drawn at two instants, with a single
marked point $\bx$ held fixed and two different material points
successively coinciding with it, would carry the idea considerably better.
As drawn, the reader can come away believing that $\bx$ is a point of each
region rather than a point of space.
\end{query}

\begin{query}{23. Figure 2.9, p.~50.}
In the rendering of the deformation of oriented surface measure, the
normals $\vek{N}$ and $\bn$ read as though they lie in the respective
tangent planes rather than perpendicular to them; the parallelograms
spanned by $d\bX,d\bY$ and $d\bx,d\by$ are drawn nearly edge-on while the
normals are drawn in the plane of the page. A perspective view in which
the parallelogram is foreshortened and the normal clearly leaves the plane
would remove the ambiguity. This figure carries the Piola--Nanson formula,
which is one of the results students most often get the geometry of wrong.
\end{query}

\begin{query}{24. Figure 2.8, p.~49.}
The three material curves intersecting at a point are shown with a
parallelepiped whose hidden edges are omitted rather than dashed, which
makes the orientation of the triad $d\bX,d\bY,d\bZ$ ambiguous, and with it
the sign of $[d\bX,d\bY,d\bZ]$ in (2.48). Since the entire discussion that
follows turns on that sign being positive, the handedness ought to be
unmistakable from the figure.
\end{query}

%=======================================================================
\section*{Summary}
%=======================================================================

\noindent
Items 1 to 9 are, I believe, straightforward misprints and are verified
above by independent computation. Item 10 is a mathematical statement that
appears stronger than the argument supports. Items 11--15 are places where
the printed equalities or terminology do not quite hold as written, though
the intended meaning is recoverable. Items 16--21 are requests for a
clarifying phrase at points where a careful reader is likely to stall, and
items 22--24 concern figures whose geometry conflicts with, or fails to
support, the accompanying text.

\end{document}
